\section{高宗时期证据链事件锚点}
以下锚点补入高宗时期财政、边疆、工程、知识与地方社会的文书链和后续节点，保持命令、条件、实施和影响之间的区分。
\subsection{1720—1729}
\eventanchor{EQ-1723-YONGZHENG-ACCESSION-DOCUMENTS}{围绕「雍正帝即位并开始重整康熙…}雍正元年（1723年），围绕「雍正帝即位并开始重整康熙晚年的财政、宗室与官僚议题」的上谕、奏折与部议在中央—地方行政链中流转。核心取证：《清世宗实录》雍正元年相关条；宫中朱批奏折、内阁大库题本与《清史稿》相关志传。这一锚点记录的是文书链：它能证明呈报、上谕、部议或照会进入机构流程，不能由此推出实际执行。来源保留：此行仅确认相关文书链和可追查的行政动作；不得由其直接推断政策有效、地方服从、民众态度或单一因果。
\eventanchor{EQ-1723-JI-ZENGYUN-CASE-DOCUMENTS}{围绕「年羹尧在西北军政中的权力…}雍正元年（1723年），围绕「年羹尧在西北军政中的权力与朝廷猜忌加剧」的上谕、奏折与部议在中央—地方行政链中流转。核心取证：《清世宗实录》雍正元年相关条；宫中朱批奏折、内阁大库题本与《清史稿》相关志传。这一锚点记录的是文书链：它能证明呈报、上谕、部议或照会进入机构流程，不能由此推出实际执行。来源保留：此行仅确认相关文书链和可追查的行政动作；不得由其直接推断政策有效、地方服从、民众态度或单一因果。
\eventanchor{EQ-1724-CATHOLIC-PROSCRIPTION-DOCUMENTS}{围绕「清廷扩大对天主教活动的限…}雍正二年（1724年），围绕「清廷扩大对天主教活动的限制，将传教、地方教民与外来人员置入治理问题」的地方呈报、案件或赈务记录将社会反应带入官府文书链。核心取证：《清世宗实录》雍正二年相关条；地方志、督抚奏折及地方档案。这一锚点记录的是文书链：它能证明呈报、上谕、部议或照会进入机构流程，不能由此推出实际执行。来源保留：此行仅确认相关文书链和可追查的行政动作；不得由其直接推断政策有效、地方服从、民众态度或单一因果。
\eventanchor{EQ-1725-NIAN-GENGYAO-PUNISHMENT-DOCUMENTS}{围绕「年羹尧被削职、赐死，其党…}雍正三年（1725年），围绕「年羹尧被削职、赐死，其党羽与相关官员受到清理」的上谕、奏折与部议在中央—地方行政链中流转。核心取证：《清世宗实录》雍正三年相关条；宫中朱批奏折、内阁大库题本与《清史稿》相关志传。这一锚点记录的是文书链：它能证明呈报、上谕、部议或照会进入机构流程，不能由此推出实际执行。来源保留：此行仅确认相关文书链和可追查的行政动作；不得由其直接推断政策有效、地方服从、民众态度或单一因果。
\eventanchor{EQ-1725-TAN-TEBI-CASE-DOCUMENTS}{围绕「隆科多失势并被幽禁，雍正…}雍正三年（1725年），围绕「隆科多失势并被幽禁，雍正初年近臣权力结构继续调整」的上谕、奏折与部议在中央—地方行政链中流转。核心取证：《清世宗实录》雍正三年相关条；宫中朱批奏折、内阁大库题本与《清史稿》相关志传。这一锚点记录的是文书链：它能证明呈报、上谕、部议或照会进入机构流程，不能由此推出实际执行。来源保留：此行仅确认相关文书链和可追查的行政动作；不得由其直接推断政策有效、地方服从、民众态度或单一因果。
\eventanchor{EQ-1726-HAOXIAN-GUIGONG-DOCUMENTS}{围绕「雍正政府推进耗羡归公与养…}雍正四年（1726年），围绕「雍正政府推进耗羡归公与养廉银安排，试图重组地方财政和官员俸给」的经费、税收、赈济或工程呈报经由户部与地方衙门处理。核心取证：《清世宗实录》雍正四年相关条；户部奏销、盐政、河工或海关档案。这一锚点记录的是文书链：它能证明呈报、上谕、部议或照会进入机构流程，不能由此推出实际执行。来源保留：此行仅确认相关文书链和可追查的行政动作；不得由其直接推断政策有效、地方服从、民众态度或单一因果。
\eventanchor{EQ-1726-Gaitu-Guiliu-SOUTHWEST-DOCUMENTS}{清廷围绕「清廷在西南继续改土归…}雍正四年（1726年），清廷围绕「清廷在西南继续改土归流，将部分土司辖地改置流官体系」处理盟旗、驻防、交通或地方呈报，相关边务文书进入中央决策链。核心取证：《清世宗实录》雍正四年相关条；军机处档、理藩院档或相关方略。这一锚点记录的是文书链：它能证明呈报、上谕、部议或照会进入机构流程，不能由此推出实际执行。来源保留：此行仅确认相关文书链和可追查的行政动作；不得由其直接推断政策有效、地方服从、民众态度或单一因果。
\eventanchor{EQ-1727-KYAKHTA-TREATY-DOCUMENTS}{围绕「《恰克图条约》确定清俄北…}雍正五年（1727年），围绕「《恰克图条约》确定清俄北部边界与贸易、使团安排」的照会、条约文本、换约或地方执行报告分别进入外交文书链。核心取证：《清世宗实录》雍正五年相关条；《筹办夷务始末》相关年卷与中外照会、条约文本。这一锚点记录的是文书链：它能证明呈报、上谕、部议或照会进入机构流程，不能由此推出实际执行。来源保留：此行仅确认相关文书链和可追查的行政动作；不得由其直接推断政策有效、地方服从、民众态度或单一因果。
\eventanchor{EQ-1727-BANNER-GENEALOGIES-DOCUMENTS}{围绕「清廷命令编修八旗谱牒，试…}雍正五年（1727年），围绕「清廷命令编修八旗谱牒，试图界定旗籍、世系与待遇资格」的地方呈报、案件或赈务记录将社会反应带入官府文书链。核心取证：《清世宗实录》雍正五年相关条；地方志、督抚奏折及地方档案。这一锚点记录的是文书链：它能证明呈报、上谕、部议或照会进入机构流程，不能由此推出实际执行。来源保留：此行仅确认相关文书链和可追查的行政动作；不得由其直接推断政策有效、地方服从、民众态度或单一因果。
\eventanchor{EQ-1728-Dzungar-EXPEDITION-DOCUMENTS}{清廷围绕「清军在西北对准噶尔展…}雍正六年（1728年），清廷围绕「清军在西北对准噶尔展开远征，后勤与高原草原路线成为军事难题」调遣军队、筹措饷械并处理战报，中央与地方的军政文书形成连续记录。核心取证：《清世宗实录》雍正六年相关条；军机处档、战事方略及地方奏报。这一锚点记录的是文书链：它能证明呈报、上谕、部议或照会进入机构流程，不能由此推出实际执行。来源保留：此行仅确认相关文书链和可追查的行政动作；不得由其直接推断政策有效、地方服从、民众态度或单一因果。
\eventanchor{EQ-1728-ZHUNGHAR-TIBET-SETTLEMENT-DOCUMENTS}{清廷围绕「清廷在准噶尔威胁下调…}雍正六年（1728年），清廷围绕「清廷在准噶尔威胁下调整驻藏与青海蒙古的军事部署」处理盟旗、驻防、交通或地方呈报，相关边务文书进入中央决策链。核心取证：《清世宗实录》雍正六年相关条；军机处档、理藩院档或相关方略。这一锚点记录的是文书链：它能证明呈报、上谕、部议或照会进入机构流程，不能由此推出实际执行。来源保留：此行仅确认相关文书链和可追查的行政动作；不得由其直接推断政策有效、地方服从、民众态度或单一因果。
\eventanchor{EQ-1729-JUNJICHU-FOUNDING-DOCUMENTS}{围绕「雍正帝设军机房，后来发展…}雍正七年（1729年），围绕「雍正帝设军机房，后来发展为军机处，处理紧急军政文书」的上谕、奏折与部议在中央—地方行政链中流转。核心取证：《清世宗实录》雍正七年相关条；宫中朱批奏折、内阁大库题本与《清史稿》相关志传。这一锚点记录的是文书链：它能证明呈报、上谕、部议或照会进入机构流程，不能由此推出实际执行。来源保留：此行仅确认相关文书链和可追查的行政动作；不得由其直接推断政策有效、地方服从、民众态度或单一因果。
\subsection{1730—1739}
\eventanchor{EQ-1730-SOUTHWEST-MIAO-GOVERNANCE-DOCUMENTS}{清廷围绕「清廷围绕苗疆设置、屯…}雍正八年（1730年），清廷围绕「清廷围绕苗疆设置、屯兵与地方纠纷加强治理」处理盟旗、驻防、交通或地方呈报，相关边务文书进入中央决策链。核心取证：《清世宗实录》雍正八年相关条；军机处档、理藩院档或相关方略。这一锚点记录的是文书链：它能证明呈报、上谕、部议或照会进入机构流程，不能由此推出实际执行。来源保留：此行仅确认相关文书链和可追查的行政动作；不得由其直接推断政策有效、地方服从、民众态度或单一因果。
\eventanchor{EQ-1731-Dzungar-WAR-REVERSE-DOCUMENTS}{清廷围绕「清军在对准噶尔战争中…}雍正九年（1731年），清廷围绕「清军在对准噶尔战争中受挫，西北远征暴露补给与指挥困难」调遣军队、筹措饷械并处理战报，中央与地方的军政文书形成连续记录。核心取证：《清世宗实录》雍正九年相关条；军机处档、战事方略及地方奏报。这一锚点记录的是文书链：它能证明呈报、上谕、部议或照会进入机构流程，不能由此推出实际执行。来源保留：此行仅确认相关文书链和可追查的行政动作；不得由其直接推断政策有效、地方服从、民众态度或单一因果。
\eventanchor{EQ-1732-KOBDO-BATTLE-DOCUMENTS}{清廷围绕「清军在和通泊等战事失…}雍正十年（1732年），清廷围绕「清军在和通泊等战事失利，西北战局一度紧张」调遣军队、筹措饷械并处理战报，中央与地方的军政文书形成连续记录。核心取证：《清世宗实录》雍正十年相关条；军机处档、战事方略及地方奏报。这一锚点记录的是文书链：它能证明呈报、上谕、部议或照会进入机构流程，不能由此推出实际执行。来源保留：此行仅确认相关文书链和可追查的行政动作；不得由其直接推断政策有效、地方服从、民众态度或单一因果。
\eventanchor{EQ-1732-AMIN-KHOJA-FLIGHT-DOCUMENTS}{清廷围绕「阿敏和卓等人在准噶尔…}雍正十年（1732年），清廷围绕「阿敏和卓等人在准噶尔压力下向清境迁移，绿洲政治进入清准竞争」处理盟旗、驻防、交通或地方呈报，相关边务文书进入中央决策链。核心取证：《清世宗实录》雍正十年相关条；军机处档、理藩院档或相关方略。这一锚点记录的是文书链：它能证明呈报、上谕、部议或照会进入机构流程，不能由此推出实际执行。来源保留：此行仅确认相关文书链和可追查的行政动作；不得由其直接推断政策有效、地方服从、民众态度或单一因果。
\eventanchor{EQ-1733-YONGZHENG-FISCAL-INSPECTION-DOCUMENTS}{围绕「雍正政府持续以奏销、亏空…}雍正十一年（1733年），围绕「雍正政府持续以奏销、亏空追查和官员考成为手段整顿地方财政」的经费、税收、赈济或工程呈报经由户部与地方衙门处理。核心取证：《清世宗实录》雍正十一年相关条；户部奏销、盐政、河工或海关档案。这一锚点记录的是文书链：它能证明呈报、上谕、部议或照会进入机构流程，不能由此推出实际执行。来源保留：此行仅确认相关文书链和可追查的行政动作；不得由其直接推断政策有效、地方服从、民众态度或单一因果。
\eventanchor{EQ-1734-CHRISTIAN-ENFORCEMENT-DOCUMENTS}{围绕「地方官继续执行限制天主教…}雍正十二年（1734年），围绕「地方官继续执行限制天主教的政策，传教士和教民的活动受地域差异影响」的地方呈报、案件或赈务记录将社会反应带入官府文书链。核心取证：《清世宗实录》雍正十二年相关条；地方志、督抚奏折及地方档案。这一锚点记录的是文书链：它能证明呈报、上谕、部议或照会进入机构流程，不能由此推出实际执行。来源保留：此行仅确认相关文书链和可追查的行政动作；不得由其直接推断政策有效、地方服从、民众态度或单一因果。
\eventanchor{EQ-1735-QIANLONG-ACCESSION-DOCUMENTS}{围绕「乾隆帝即位，雍正时期形成…}雍正十三年（1735年），围绕「乾隆帝即位，雍正时期形成的财政与军机决策机制进入新朝」的上谕、奏折与部议在中央—地方行政链中流转。核心取证：《清世宗实录》雍正十三年相关条；宫中朱批奏折、内阁大库题本与《清史稿》相关志传。这一锚点记录的是文书链：它能证明呈报、上谕、部议或照会进入机构流程，不能由此推出实际执行。来源保留：此行仅确认相关文书链和可追查的行政动作；不得由其直接推断政策有效、地方服从、民众态度或单一因果。
\eventanchor{EQ-1735-MIAO-UPRISING-RONGJIANG-DOCUMENTS}{清廷围绕「贵州荣江等地苗侗社区…}雍正十三年（1735年），清廷围绕「贵州荣江等地苗侗社区反抗清军和地方治理，镇压造成严重暴力」调遣军队、筹措饷械并处理战报，中央与地方的军政文书形成连续记录。核心取证：《清世宗实录》雍正十三年相关条；军机处档、战事方略及地方奏报。这一锚点记录的是文书链：它能证明呈报、上谕、部议或照会进入机构流程，不能由此推出实际执行。来源保留：此行仅确认相关文书链和可追查的行政动作；不得由其直接推断政策有效、地方服从、民众态度或单一因果。
\eventanchor{EQ-1736-AMNESTY-AND-RECTIFICATION-DOCUMENTS}{围绕「乾隆初年颁布恩诏并调整雍…}乾隆元年（1736年），围绕「乾隆初年颁布恩诏并调整雍正后期政治案件的处置」的上谕、奏折与部议在中央—地方行政链中流转。核心取证：《清高宗实录》乾隆元年相关条；宫中朱批奏折、内阁大库题本与《清史稿》相关志传。这一锚点记录的是文书链：它能证明呈报、上谕、部议或照会进入机构流程，不能由此推出实际执行。来源保留：此行仅确认相关文书链和可追查的行政动作；不得由其直接推断政策有效、地方服从、民众态度或单一因果。
\eventanchor{EQ-1737-Dzungar-NEGOTIATION-DOCUMENTS}{围绕「清廷在战争压力下继续处理…}乾隆二年（1737年），围绕「清廷在战争压力下继续处理与准噶尔的使节、盟旗和军事情报」的照会、条约文本、换约或地方执行报告分别进入外交文书链。核心取证：《清高宗实录》乾隆二年相关条；《筹办夷务始末》相关年卷与中外照会、条约文本。这一锚点记录的是文书链：它能证明呈报、上谕、部议或照会进入机构流程，不能由此推出实际执行。来源保留：此行仅确认相关文书链和可追查的行政动作；不得由其直接推断政策有效、地方服从、民众态度或单一因果。
\eventanchor{EQ-1738-FIRST-JINCHUAN-WAR-DOCUMENTS}{清廷围绕「清军开始大规模干预大…}乾隆三年（1738年），清廷围绕「清军开始大规模干预大金川，山地战与土司关系限制军事推进」调遣军队、筹措饷械并处理战报，中央与地方的军政文书形成连续记录。核心取证：《清高宗实录》乾隆三年相关条；军机处档、战事方略及地方奏报。这一锚点记录的是文书链：它能证明呈报、上谕、部议或照会进入机构流程，不能由此推出实际执行。来源保留：此行仅确认相关文书链和可追查的行政动作；不得由其直接推断政策有效、地方服从、民众态度或单一因果。
\eventanchor{EQ-1739-JINCHUAN-SUPPLY-DOCUMENTS}{围绕「大金川战事的粮运、道路和…}乾隆四年（1739年），围绕「大金川战事的粮运、道路和军费成为中央持续关注的难题」的经费、税收、赈济或工程呈报经由户部与地方衙门处理。核心取证：《清高宗实录》乾隆四年相关条；户部奏销、盐政、河工或海关档案。这一锚点记录的是文书链：它能证明呈报、上谕、部议或照会进入机构流程，不能由此推出实际执行。来源保留：此行仅确认相关文书链和可追查的行政动作；不得由其直接推断政策有效、地方服从、民众态度或单一因果。
\subsection{1740—1749}
\eventanchor{EQ-1740-NORTH-CHINA-FAMINE-DOCUMENTS}{围绕「华北旱灾和粮价问题促使地…}乾隆五年（1740年），围绕「华北旱灾和粮价问题促使地方奏报与赈济措施增加」的地方呈报、案件或赈务记录将社会反应带入官府文书链。核心取证：《清高宗实录》乾隆五年相关条；地方志、督抚奏折及地方档案。这一锚点记录的是文书链：它能证明呈报、上谕、部议或照会进入机构流程，不能由此推出实际执行。来源保留：此行仅确认相关文书链和可追查的行政动作；不得由其直接推断政策有效、地方服从、民众态度或单一因果。
\eventanchor{EQ-1741-MIAO-FRONTIER-CAMPAIGN-DOCUMENTS}{清廷围绕「清廷在西南继续以驻兵…}乾隆六年（1741年），清廷围绕「清廷在西南继续以驻兵、设治和军事行动处理苗疆冲突」处理盟旗、驻防、交通或地方呈报，相关边务文书进入中央决策链。核心取证：《清高宗实录》乾隆六年相关条；军机处档、理藩院档或相关方略。这一锚点记录的是文书链：它能证明呈报、上谕、部议或照会进入机构流程，不能由此推出实际执行。来源保留：此行仅确认相关文书链和可追查的行政动作；不得由其直接推断政策有效、地方服从、民众态度或单一因果。
\eventanchor{EQ-1742-BANNER-EXIT-ORDER-DOCUMENTS}{围绕「清廷允许部分汉军旗人出旗…}乾隆七年（1742年），围绕「清廷允许部分汉军旗人出旗，调整旗籍与生计安排」的地方呈报、案件或赈务记录将社会反应带入官府文书链。核心取证：《清高宗实录》乾隆七年相关条；地方志、督抚奏折及地方档案。这一锚点记录的是文书链：它能证明呈报、上谕、部议或照会进入机构流程，不能由此推出实际执行。来源保留：此行仅确认相关文书链和可追查的行政动作；不得由其直接推断政策有效、地方服从、民众态度或单一因果。
\eventanchor{EQ-1743-YELLOW-RIVER-WORKS-DOCUMENTS}{围绕「河工官员就黄河堤防、河道…}乾隆八年（1743年），围绕「河工官员就黄河堤防、河道和经费持续奏报并实施工程」的经费、税收、赈济或工程呈报经由户部与地方衙门处理。核心取证：《清高宗实录》乾隆八年相关条；户部奏销、盐政、河工或海关档案。这一锚点记录的是文书链：它能证明呈报、上谕、部议或照会进入机构流程，不能由此推出实际执行。来源保留：此行仅确认相关文书链和可追查的行政动作；不得由其直接推断政策有效、地方服从、民众态度或单一因果。
\eventanchor{EQ-1744-JESUIT-CALENDAR-SERVICE-DOCUMENTS}{围绕「宫廷继续利用部分西洋传教…}乾隆九年（1744年），围绕「宫廷继续利用部分西洋传教士的历算与技术服务，同时维持宗教限制」的禁令、编纂、教育或传播活动留下中央和地方文本记录。核心取证：《清高宗实录》乾隆九年相关条；内阁档、文集或报刊相关记录。这一锚点记录的是文书链：它能证明呈报、上谕、部议或照会进入机构流程，不能由此推出实际执行。来源保留：此行仅确认相关文书链和可追查的行政动作；不得由其直接推断政策有效、地方服从、民众态度或单一因果。
\eventanchor{EQ-1745-JINCHUAN-NEGOTIATION-DOCUMENTS}{清廷围绕「清廷在大金川战局中尝…}乾隆十年（1745年），清廷围绕「清廷在大金川战局中尝试通过招抚、盟约和军事压力重组地方关系」处理盟旗、驻防、交通或地方呈报，相关边务文书进入中央决策链。核心取证：《清高宗实录》乾隆十年相关条；军机处档、理藩院档或相关方略。这一锚点记录的是文书链：它能证明呈报、上谕、部议或照会进入机构流程，不能由此推出实际执行。来源保留：此行仅确认相关文书链和可追查的行政动作；不得由其直接推断政策有效、地方服从、民众态度或单一因果。
\eventanchor{EQ-1746-FIRST-JINCHUAN-SETTLEMENT-DOCUMENTS}{清廷围绕「第一次金川战争以和议…}乾隆十一年（1746年），清廷围绕「第一次金川战争以和议和行政安排告一段落」调遣军队、筹措饷械并处理战报，中央与地方的军政文书形成连续记录。核心取证：《清高宗实录》乾隆十一年相关条；军机处档、战事方略及地方奏报。这一锚点记录的是文书链：它能证明呈报、上谕、部议或照会进入机构流程，不能由此推出实际执行。来源保留：此行仅确认相关文书链和可追查的行政动作；不得由其直接推断政策有效、地方服从、民众态度或单一因果。
\eventanchor{EQ-1747-SOUTHWEST-MINING-REGULATION-DOCUMENTS}{围绕「清廷和地方官围绕银铜矿开…}乾隆十二年（1747年），围绕「清廷和地方官围绕银铜矿开采、税课与矿工流动加强管理」的经费、税收、赈济或工程呈报经由户部与地方衙门处理。核心取证：《清高宗实录》乾隆十二年相关条；户部奏销、盐政、河工或海关档案。这一锚点记录的是文书链：它能证明呈报、上谕、部议或照会进入机构流程，不能由此推出实际执行。来源保留：此行仅确认相关文书链和可追查的行政动作；不得由其直接推断政策有效、地方服从、民众态度或单一因果。
\eventanchor{EQ-1748-SONG-AND-CASE-DOCUMENTS}{围绕「乾隆朝通过大案处分官员，…}乾隆十三年（1748年），围绕「乾隆朝通过大案处分官员，强化对地方文书、钱粮和吏治的监督」的上谕、奏折与部议在中央—地方行政链中流转。核心取证：《清高宗实录》乾隆十三年相关条；宫中朱批奏折、内阁大库题本与《清史稿》相关志传。这一锚点记录的是文书链：它能证明呈报、上谕、部议或照会进入机构流程，不能由此推出实际执行。来源保留：此行仅确认相关文书链和可追查的行政动作；不得由其直接推断政策有效、地方服从、民众态度或单一因果。
\eventanchor{EQ-1749-JIANGNAN-TAX-REVIEW-DOCUMENTS}{围绕「江南钱粮、漕运和士绅关系…}乾隆十四年（1749年），围绕「江南钱粮、漕运和士绅关系继续在户部与地方奏折中被讨论」的经费、税收、赈济或工程呈报经由户部与地方衙门处理。核心取证：《清高宗实录》乾隆十四年相关条；户部奏销、盐政、河工或海关档案。这一锚点记录的是文书链：它能证明呈报、上谕、部议或照会进入机构流程，不能由此推出实际执行。来源保留：此行仅确认相关文书链和可追查的行政动作；不得由其直接推断政策有效、地方服从、民众态度或单一因果。
\subsection{1750—1759}
\eventanchor{EQ-1750-LHASA-UPRISING-DOCUMENTS}{清廷围绕「拉萨发生政治冲突，驻…}乾隆十五年（1750年），清廷围绕「拉萨发生政治冲突，驻藏大臣和清军介入处置」处理盟旗、驻防、交通或地方呈报，相关边务文书进入中央决策链。核心取证：《清高宗实录》乾隆十五年相关条；军机处档、理藩院档或相关方略。这一锚点记录的是文书链：它能证明呈报、上谕、部议或照会进入机构流程，不能由此推出实际执行。来源保留：此行仅确认相关文书链和可追查的行政动作；不得由其直接推断政策有效、地方服从、民众态度或单一因果。
\eventanchor{EQ-1751-TIBET-ORDINANCE-DOCUMENTS}{清廷围绕「清廷在拉萨事变后重整…}乾隆十六年（1751年），清廷围绕「清廷在拉萨事变后重整驻藏大臣和噶厦相关制度安排」处理盟旗、驻防、交通或地方呈报，相关边务文书进入中央决策链。核心取证：《清高宗实录》乾隆十六年相关条；军机处档、理藩院档或相关方略。这一锚点记录的是文书链：它能证明呈报、上谕、部议或照会进入机构流程，不能由此推出实际执行。来源保留：此行仅确认相关文书链和可追查的行政动作；不得由其直接推断政策有效、地方服从、民众态度或单一因果。
\eventanchor{EQ-1752-TIBETAN-ADMINISTRATION-DOCUMENTS}{清廷围绕「驻藏系统与青海、川藏…}乾隆十七年（1752年），清廷围绕「驻藏系统与青海、川藏交通的军政协调继续调整」处理盟旗、驻防、交通或地方呈报，相关边务文书进入中央决策链。核心取证：《清高宗实录》乾隆十七年相关条；军机处档、理藩院档或相关方略。这一锚点记录的是文书链：它能证明呈报、上谕、部议或照会进入机构流程，不能由此推出实际执行。来源保留：此行仅确认相关文书链和可追查的行政动作；不得由其直接推断政策有效、地方服从、民众态度或单一因果。
\eventanchor{EQ-1753-SALT-ADMINISTRATION-DOCUMENTS}{围绕「盐政、引岸和课额问题在中…}乾隆十八年（1753年），围绕「盐政、引岸和课额问题在中央与产销地区持续受到整顿」的经费、税收、赈济或工程呈报经由户部与地方衙门处理。核心取证：《清高宗实录》乾隆十八年相关条；户部奏销、盐政、河工或海关档案。这一锚点记录的是文书链：它能证明呈报、上谕、部议或照会进入机构流程，不能由此推出实际执行。来源保留：此行仅确认相关文书链和可追查的行政动作；不得由其直接推断政策有效、地方服从、民众态度或单一因果。
\eventanchor{EQ-1754-AMURSANA-DEFECTION-DOCUMENTS}{清廷围绕「阿睦尔撒纳等准噶尔贵…}乾隆十九年（1754年），清廷围绕「阿睦尔撒纳等准噶尔贵族向清廷归附，改变西北战争的联盟格局」处理盟旗、驻防、交通或地方呈报，相关边务文书进入中央决策链。核心取证：《清高宗实录》乾隆十九年相关条；军机处档、理藩院档或相关方略。这一锚点记录的是文书链：它能证明呈报、上谕、部议或照会进入机构流程，不能由此推出实际执行。来源保留：此行仅确认相关文书链和可追查的行政动作；不得由其直接推断政策有效、地方服从、民众态度或单一因果。
\eventanchor{EQ-1755-ILI-EXPEDITION-DOCUMENTS}{清廷围绕「清军分路进入伊犁，击…}乾隆二十年（1755年），清廷围绕「清军分路进入伊犁，击溃达瓦齐政权并占领准噶尔核心地区」调遣军队、筹措饷械并处理战报，中央与地方的军政文书形成连续记录。核心取证：《清高宗实录》乾隆二十年相关条；军机处档、战事方略及地方奏报。这一锚点记录的是文书链：它能证明呈报、上谕、部议或照会进入机构流程，不能由此推出实际执行。来源保留：此行仅确认相关文书链和可追查的行政动作；不得由其直接推断政策有效、地方服从、民众态度或单一因果。
\eventanchor{EQ-1755-AMURSANA-REVOLT-DOCUMENTS}{清廷围绕「阿睦尔撒纳反清，准噶…}乾隆二十年（1755年），清廷围绕「阿睦尔撒纳反清，准噶尔战后秩序重新陷入军事冲突」调遣军队、筹措饷械并处理战报，中央与地方的军政文书形成连续记录。核心取证：《清高宗实录》乾隆二十年相关条；军机处档、战事方略及地方奏报。这一锚点记录的是文书链：它能证明呈报、上谕、部议或照会进入机构流程，不能由此推出实际执行。来源保留：此行仅确认相关文书链和可追查的行政动作；不得由其直接推断政策有效、地方服从、民众态度或单一因果。
\eventanchor{EQ-1756-BANNER-CIVILIAN-REGISTRATION-DOCUMENTS}{围绕「清廷命令部分八旗副户口改…}乾隆二十一年（1756年），围绕「清廷命令部分八旗副户口改归民籍，继续缓解旗籍供养压力」的地方呈报、案件或赈务记录将社会反应带入官府文书链。核心取证：《清高宗实录》乾隆二十一年相关条；地方志、督抚奏折及地方档案。这一锚点记录的是文书链：它能证明呈报、上谕、部议或照会进入机构流程，不能由此推出实际执行。来源保留：此行仅确认相关文书链和可追查的行政动作；不得由其直接推断政策有效、地方服从、民众态度或单一因果。
\eventanchor{EQ-1757-CANTON-TRADE-RESTRICTION-DOCUMENTS}{围绕「清廷将西洋贸易集中于广州…}乾隆二十二年（1757年），围绕「清廷将西洋贸易集中于广州口岸，形成后来所谓“一口通商”的制度框架」的照会、条约文本、换约或地方执行报告分别进入外交文书链。核心取证：《清高宗实录》乾隆二十二年相关条；《筹办夷务始末》相关年卷与中外照会、条约文本。这一锚点记录的是文书链：它能证明呈报、上谕、部议或照会进入机构流程，不能由此推出实际执行。来源保留：此行仅确认相关文书链和可追查的行政动作；不得由其直接推断政策有效、地方服从、民众态度或单一因果。
\eventanchor{EQ-1757-AMURSANA-DEATH-DOCUMENTS}{清廷围绕「阿睦尔撒纳逃入俄境后…}乾隆二十二年（1757年），清廷围绕「阿睦尔撒纳逃入俄境后死亡，清准战争的关键对手退出」处理盟旗、驻防、交通或地方呈报，相关边务文书进入中央决策链。核心取证：《清高宗实录》乾隆二十二年相关条；军机处档、理藩院档或相关方略。这一锚点记录的是文书链：它能证明呈报、上谕、部议或照会进入机构流程，不能由此推出实际执行。来源保留：此行仅确认相关文书链和可追查的行政动作；不得由其直接推断政策有效、地方服从、民众态度或单一因果。
\eventanchor{EQ-1758-KHOJA-REBELLION-DOCUMENTS}{清廷围绕「大小和卓反清，清军由…}乾隆二十三年（1758年），清廷围绕「大小和卓反清，清军由北路向天山南路推进」调遣军队、筹措饷械并处理战报，中央与地方的军政文书形成连续记录。核心取证：《清高宗实录》乾隆二十三年相关条；军机处档、战事方略及地方奏报。这一锚点记录的是文书链：它能证明呈报、上谕、部议或照会进入机构流程，不能由此推出实际执行。来源保留：此行仅确认相关文书链和可追查的行政动作；不得由其直接推断政策有效、地方服从、民众态度或单一因果。
\eventanchor{EQ-1759-KASHGAR-YARKAND-CAMPAIGN-DOCUMENTS}{清廷围绕「清军攻入喀什噶尔、叶…}乾隆二十四年（1759年），清廷围绕「清军攻入喀什噶尔、叶尔羌等地，和卓政权覆亡」调遣军队、筹措饷械并处理战报，中央与地方的军政文书形成连续记录。核心取证：《清高宗实录》乾隆二十四年相关条；军机处档、战事方略及地方奏报。这一锚点记录的是文书链：它能证明呈报、上谕、部议或照会进入机构流程，不能由此推出实际执行。来源保留：此行仅确认相关文书链和可追查的行政动作；不得由其直接推断政策有效、地方服从、民众态度或单一因果。
\subsection{1760—1769}
\eventanchor{EQ-1760-ILI-GENERAL-DOCUMENTS}{清廷围绕「清廷设伊犁将军并展开…}乾隆二十五年（1760年），清廷围绕「清廷设伊犁将军并展开驻防、屯田与移民安置」处理盟旗、驻防、交通或地方呈报，相关边务文书进入中央决策链。核心取证：《清高宗实录》乾隆二十五年相关条；军机处档、理藩院档或相关方略。这一锚点记录的是文书链：它能证明呈报、上谕、部议或照会进入机构流程，不能由此推出实际执行。来源保留：此行仅确认相关文书链和可追查的行政动作；不得由其直接推断政策有效、地方服从、民众态度或单一因果。
\eventanchor{EQ-1761-XINJIANG-GARRISONS-DOCUMENTS}{清廷围绕「清廷在新疆配置满汉蒙…}乾隆二十六年（1761年），清廷围绕「清廷在新疆配置满汉蒙军队和绿营驻防，建立军政节点」处理盟旗、驻防、交通或地方呈报，相关边务文书进入中央决策链。核心取证：《清高宗实录》乾隆二十六年相关条；军机处档、理藩院档或相关方略。这一锚点记录的是文书链：它能证明呈报、上谕、部议或照会进入机构流程，不能由此推出实际执行。来源保留：此行仅确认相关文书链和可追查的行政动作；不得由其直接推断政策有效、地方服从、民众态度或单一因果。
\eventanchor{EQ-1762-URUMQI-FOUNDATION-DOCUMENTS}{清廷围绕「迪化等北疆城市的军政…}乾隆二十七年（1762年），清廷围绕「迪化等北疆城市的军政建设加快，成为屯田和交通枢纽」处理盟旗、驻防、交通或地方呈报，相关边务文书进入中央决策链。核心取证：《清高宗实录》乾隆二十七年相关条；军机处档、理藩院档或相关方略。这一锚点记录的是文书链：它能证明呈报、上谕、部议或照会进入机构流程，不能由此推出实际执行。来源保留：此行仅确认相关文书链和可追查的行政动作；不得由其直接推断政策有效、地方服从、民众态度或单一因果。
\eventanchor{EQ-1763-XINJIANG-SETTLEMENT-DOCUMENTS}{围绕「清廷继续在新疆安置屯田人…}乾隆二十八年（1763年），围绕「清廷继续在新疆安置屯田人口、军户和商业网络」的地方呈报、案件或赈务记录将社会反应带入官府文书链。核心取证：《清高宗实录》乾隆二十八年相关条；地方志、督抚奏折及地方档案。这一锚点记录的是文书链：它能证明呈报、上谕、部议或照会进入机构流程，不能由此推出实际执行。来源保留：此行仅确认相关文书链和可追查的行政动作；不得由其直接推断政策有效、地方服从、民众态度或单一因果。
\eventanchor{EQ-1764-MANCHU-BANNER-RESETTLEMENT-DOCUMENTS}{围绕「部分八旗兵丁被调往西北驻…}乾隆二十九年（1764年），围绕「部分八旗兵丁被调往西北驻防，旗籍家庭的迁移与供给成为边务问题」的地方呈报、案件或赈务记录将社会反应带入官府文书链。核心取证：《清高宗实录》乾隆二十九年相关条；地方志、督抚奏折及地方档案。这一锚点记录的是文书链：它能证明呈报、上谕、部议或照会进入机构流程，不能由此推出实际执行。来源保留：此行仅确认相关文书链和可追查的行政动作；不得由其直接推断政策有效、地方服从、民众态度或单一因果。
\eventanchor{EQ-1765-USH-REBELLION-DOCUMENTS}{清廷围绕「乌什爆发反清事件，清…}乾隆三十年（1765年），清廷围绕「乌什爆发反清事件，清军镇压并重组当地治理」调遣军队、筹措饷械并处理战报，中央与地方的军政文书形成连续记录。核心取证：《清高宗实录》乾隆三十年相关条；军机处档、战事方略及地方奏报。这一锚点记录的是文书链：它能证明呈报、上谕、部议或照会进入机构流程，不能由此推出实际执行。来源保留：此行仅确认相关文书链和可追查的行政动作；不得由其直接推断政策有效、地方服从、民众态度或单一因果。
\eventanchor{EQ-1766-BURMA-WAR-BEGIN-DOCUMENTS}{清廷围绕「清军对缅甸发动远征，…}乾隆三十一年（1766年），清廷围绕「清军对缅甸发动远征，西南边境战争展开」调遣军队、筹措饷械并处理战报，中央与地方的军政文书形成连续记录。核心取证：《清高宗实录》乾隆三十一年相关条；军机处档、战事方略及地方奏报。这一锚点记录的是文书链：它能证明呈报、上谕、部议或照会进入机构流程，不能由此推出实际执行。来源保留：此行仅确认相关文书链和可追查的行政动作；不得由其直接推断政策有效、地方服从、民众态度或单一因果。
\eventanchor{EQ-1767-BURMA-CAMPAIGN-DOCUMENTS}{清廷围绕「清缅战争持续，清军在…}乾隆三十二年（1767年），清廷围绕「清缅战争持续，清军在边境和山地路线中遭遇重大困难」调遣军队、筹措饷械并处理战报，中央与地方的军政文书形成连续记录。核心取证：《清高宗实录》乾隆三十二年相关条；军机处档、战事方略及地方奏报。这一锚点记录的是文书链：它能证明呈报、上谕、部议或照会进入机构流程，不能由此推出实际执行。来源保留：此行仅确认相关文书链和可追查的行政动作；不得由其直接推断政策有效、地方服从、民众态度或单一因果。
\eventanchor{EQ-1768-LITERARY-INQUISITION-DOCUMENTS}{围绕「乾隆朝以文字狱案件强化对…}乾隆三十三年（1768年），围绕「乾隆朝以文字狱案件强化对著述、地方士人和政治语言的控制」的禁令、编纂、教育或传播活动留下中央和地方文本记录。核心取证：《清高宗实录》乾隆三十三年相关条；内阁档、文集或报刊相关记录。这一锚点记录的是文书链：它能证明呈报、上谕、部议或照会进入机构流程，不能由此推出实际执行。来源保留：此行仅确认相关文书链和可追查的行政动作；不得由其直接推断政策有效、地方服从、民众态度或单一因果。
\eventanchor{EQ-1768-SECOND-JINCHUAN-WAR-DOCUMENTS}{清廷围绕「第二次金川战争开始，…}乾隆三十三年（1768年），清廷围绕「第二次金川战争开始，清廷投入更大兵力和资源」调遣军队、筹措饷械并处理战报，中央与地方的军政文书形成连续记录。核心取证：《清高宗实录》乾隆三十三年相关条；军机处档、战事方略及地方奏报。这一锚点记录的是文书链：它能证明呈报、上谕、部议或照会进入机构流程，不能由此推出实际执行。来源保留：此行仅确认相关文书链和可追查的行政动作；不得由其直接推断政策有效、地方服从、民众态度或单一因果。
\eventanchor{EQ-1769-BURMA-SETTLEMENT-DOCUMENTS}{围绕「清缅战争在反复交战后趋向…}乾隆三十四年（1769年），围绕「清缅战争在反复交战后趋向停战和使节往来安排」的照会、条约文本、换约或地方执行报告分别进入外交文书链。核心取证：《清高宗实录》乾隆三十四年相关条；《筹办夷务始末》相关年卷与中外照会、条约文本。这一锚点记录的是文书链：它能证明呈报、上谕、部议或照会进入机构流程，不能由此推出实际执行。来源保留：此行仅确认相关文书链和可追查的行政动作；不得由其直接推断政策有效、地方服从、民众态度或单一因果。
\eventanchor{EQ-1769-JINCHUAN-ESCALATION-DOCUMENTS}{清廷围绕「金川战场的道路、炮兵…}乾隆三十四年（1769年），清廷围绕「金川战场的道路、炮兵、粮运与将领调度不断升级」调遣军队、筹措饷械并处理战报，中央与地方的军政文书形成连续记录。核心取证：《清高宗实录》乾隆三十四年相关条；军机处档、战事方略及地方奏报。这一锚点记录的是文书链：它能证明呈报、上谕、部议或照会进入机构流程，不能由此推出实际执行。来源保留：此行仅确认相关文书链和可追查的行政动作；不得由其直接推断政策有效、地方服从、民众态度或单一因果。
\subsection{1770—1779}
\eventanchor{EQ-1770-SICHUAN-SUPPLY-BURDEN-DOCUMENTS}{围绕「金川军需加重四川及邻省的…}乾隆三十五年（1770年），围绕「金川军需加重四川及邻省的转运、民夫和财政负担」的经费、税收、赈济或工程呈报经由户部与地方衙门处理。核心取证：《清高宗实录》乾隆三十五年相关条；户部奏销、盐政、河工或海关档案。这一锚点记录的是文书链：它能证明呈报、上谕、部议或照会进入机构流程，不能由此推出实际执行。来源保留：此行仅确认相关文书链和可追查的行政动作；不得由其直接推断政策有效、地方服从、民众态度或单一因果。
\eventanchor{EQ-1771-TORGHUT-RETURN-DOCUMENTS}{清廷围绕「土尔扈特部自伏尔加回…}乾隆三十六年（1771年），清廷围绕「土尔扈特部自伏尔加回归，清廷在伊犁等地安置部众」处理盟旗、驻防、交通或地方呈报，相关边务文书进入中央决策链。核心取证：《清高宗实录》乾隆三十六年相关条；军机处档、理藩院档或相关方略。这一锚点记录的是文书链：它能证明呈报、上谕、部议或照会进入机构流程，不能由此推出实际执行。来源保留：此行仅确认相关文书链和可追查的行政动作；不得由其直接推断政策有效、地方服从、民众态度或单一因果。
\eventanchor{EQ-1772-SIKU-QUANSHU-ORDER-DOCUMENTS}{围绕「乾隆帝下令征集、编纂《四…}乾隆三十七年（1772年），围绕「乾隆帝下令征集、编纂《四库全书》，建立全国性的图书征集机制」的禁令、编纂、教育或传播活动留下中央和地方文本记录。核心取证：《清高宗实录》乾隆三十七年相关条；内阁档、文集或报刊相关记录。这一锚点记录的是文书链：它能证明呈报、上谕、部议或照会进入机构流程，不能由此推出实际执行。来源保留：此行仅确认相关文书链和可追查的行政动作；不得由其直接推断政策有效、地方服从、民众态度或单一因果。
\eventanchor{EQ-1773-SICHUAN-RECRUITMENT-DOCUMENTS}{清廷围绕「金川战事促使清廷继续…}乾隆三十八年（1773年），清廷围绕「金川战事促使清廷继续在四川及邻省征调兵丁和夫役」调遣军队、筹措饷械并处理战报，中央与地方的军政文书形成连续记录。核心取证：《清高宗实录》乾隆三十八年相关条；军机处档、战事方略及地方奏报。这一锚点记录的是文书链：它能证明呈报、上谕、部议或照会进入机构流程，不能由此推出实际执行。来源保留：此行仅确认相关文书链和可追查的行政动作；不得由其直接推断政策有效、地方服从、民众态度或单一因果。
\eventanchor{EQ-1774-WANG-LUN-UPRISING-DOCUMENTS}{围绕「王伦在山东起事，清廷迅速…}乾隆三十九年（1774年），围绕「王伦在山东起事，清廷迅速调兵镇压」的地方呈报、案件或赈务记录将社会反应带入官府文书链。核心取证：《清高宗实录》乾隆三十九年相关条；地方志、督抚奏折及地方档案。这一锚点记录的是文书链：它能证明呈报、上谕、部议或照会进入机构流程，不能由此推出实际执行。来源保留：此行仅确认相关文书链和可追查的行政动作；不得由其直接推断政策有效、地方服从、民众态度或单一因果。
\eventanchor{EQ-1775-SIKU-COLLECTION-DOCUMENTS}{围绕「《四库全书》征书与审查在…}乾隆四十年（1775年），围绕「《四库全书》征书与审查在各省展开，地方藏书和文人网络受到影响」的禁令、编纂、教育或传播活动留下中央和地方文本记录。核心取证：《清高宗实录》乾隆四十年相关条；内阁档、文集或报刊相关记录。这一锚点记录的是文书链：它能证明呈报、上谕、部议或照会进入机构流程，不能由此推出实际执行。来源保留：此行仅确认相关文书链和可追查的行政动作；不得由其直接推断政策有效、地方服从、民众态度或单一因果。
\eventanchor{EQ-1776-JINCHUAN-CONQUEST-DOCUMENTS}{清廷围绕「清军攻克金川，第二次…}乾隆四十一年（1776年），清廷围绕「清军攻克金川，第二次金川战争结束」调遣军队、筹措饷械并处理战报，中央与地方的军政文书形成连续记录。核心取证：《清高宗实录》乾隆四十一年相关条；军机处档、战事方略及地方奏报。这一锚点记录的是文书链：它能证明呈报、上谕、部议或照会进入机构流程，不能由此推出实际执行。来源保留：此行仅确认相关文书链和可追查的行政动作；不得由其直接推断政策有效、地方服从、民众态度或单一因果。
\eventanchor{EQ-1777-JINCHUAN-ADMINISTRATIVE-RESET-DOCUMENTS}{清廷围绕「清廷在金川地区推行改…}乾隆四十二年（1777年），清廷围绕「清廷在金川地区推行改土归流、设治和驻防安排」处理盟旗、驻防、交通或地方呈报，相关边务文书进入中央决策链。核心取证：《清高宗实录》乾隆四十二年相关条；军机处档、理藩院档或相关方略。这一锚点记录的是文书链：它能证明呈报、上谕、部议或照会进入机构流程，不能由此推出实际执行。来源保留：此行仅确认相关文书链和可追查的行政动作；不得由其直接推断政策有效、地方服从、民众态度或单一因果。
\eventanchor{EQ-1778-SICHUAN-TAX-RECOVERY-DOCUMENTS}{围绕「金川战后四川财政和民力恢…}乾隆四十三年（1778年），围绕「金川战后四川财政和民力恢复成为地方官奏报的重要议题」的经费、税收、赈济或工程呈报经由户部与地方衙门处理。核心取证：《清高宗实录》乾隆四十三年相关条；户部奏销、盐政、河工或海关档案。这一锚点记录的是文书链：它能证明呈报、上谕、部议或照会进入机构流程，不能由此推出实际执行。来源保留：此行仅确认相关文书链和可追查的行政动作；不得由其直接推断政策有效、地方服从、民众态度或单一因果。
\eventanchor{EQ-1779-VIETNAM-INTERVENTION-DOCUMENTS}{清廷围绕「清廷介入安南政局并派…}乾隆四十四年（1779年），清廷围绕「清廷介入安南政局并派兵南下，孙士毅军一度进入升龙」调遣军队、筹措饷械并处理战报，中央与地方的军政文书形成连续记录。核心取证：《清高宗实录》乾隆四十四年相关条；军机处档、战事方略及地方奏报。这一锚点记录的是文书链：它能证明呈报、上谕、部议或照会进入机构流程，不能由此推出实际执行。来源保留：此行仅确认相关文书链和可追查的行政动作；不得由其直接推断政策有效、地方服从、民众态度或单一因果。
\subsection{1780—1789}
\eventanchor{EQ-1780-VIETNAM-WITHDRAWAL-DOCUMENTS}{围绕「清军撤离安南，双方通过使…}乾隆四十五年（1780年），围绕「清军撤离安南，双方通过使节和册封语言恢复关系」的照会、条约文本、换约或地方执行报告分别进入外交文书链。核心取证：《清高宗实录》乾隆四十五年相关条；《筹办夷务始末》相关年卷与中外照会、条约文本。这一锚点记录的是文书链：它能证明呈报、上谕、部议或照会进入机构流程，不能由此推出实际执行。来源保留：此行仅确认相关文书链和可追查的行政动作；不得由其直接推断政策有效、地方服从、民众态度或单一因果。
\eventanchor{EQ-1781-GANSU-MUSLIM-UPRISING-DOCUMENTS}{清廷围绕「甘肃爆发回民起事，清…}乾隆四十六年（1781年），清廷围绕「甘肃爆发回民起事，清军实施镇压与后续治理」调遣军队、筹措饷械并处理战报，中央与地方的军政文书形成连续记录。核心取证：《清高宗实录》乾隆四十六年相关条；军机处档、战事方略及地方奏报。这一锚点记录的是文书链：它能证明呈报、上谕、部议或照会进入机构流程，不能由此推出实际执行。来源保留：此行仅确认相关文书链和可追查的行政动作；不得由其直接推断政策有效、地方服从、民众态度或单一因果。
\eventanchor{EQ-1782-GANSU-POSTWAR-GOVERNANCE-DOCUMENTS}{围绕「清廷在甘肃起事后调整驻兵…}乾隆四十七年（1782年），围绕「清廷在甘肃起事后调整驻兵、司法和地方控制措施」的地方呈报、案件或赈务记录将社会反应带入官府文书链。核心取证：《清高宗实录》乾隆四十七年相关条；地方志、督抚奏折及地方档案。这一锚点记录的是文书链：它能证明呈报、上谕、部议或照会进入机构流程，不能由此推出实际执行。来源保留：此行仅确认相关文书链和可追查的行政动作；不得由其直接推断政策有效、地方服从、民众态度或单一因果。
\eventanchor{EQ-1783-SALT-MONOPOLY-REVIEW-DOCUMENTS}{围绕「清廷继续讨论盐政弊端、商…}乾隆四十八年（1783年），围绕「清廷继续讨论盐政弊端、商人负担和课额征解」的经费、税收、赈济或工程呈报经由户部与地方衙门处理。核心取证：《清高宗实录》乾隆四十八年相关条；户部奏销、盐政、河工或海关档案。这一锚点记录的是文书链：它能证明呈报、上谕、部议或照会进入机构流程，不能由此推出实际执行。来源保留：此行仅确认相关文书链和可追查的行政动作；不得由其直接推断政策有效、地方服从、民众态度或单一因果。
\eventanchor{EQ-1784-TAIWAN-UPRISING-PRELUDE}{台湾地方械斗、赋役与移民社会紧…}乾隆四十九年（1784年），台湾地方械斗、赋役与移民社会紧张为大规模反抗累积条件。核心取证：《清高宗实录》乾隆四十九年相关条；地方志、督抚奏折及地方档案。这一锚点界定可回查的前史条件；条件、传闻或信息累积不应被写成已经证实的单一直接原因。来源保留：以乾隆四十九年（1784年）的同期文书为定位起点；官修记录、奏报或外交文本服务特定机构立场，具体日期、规模、地方执行与对手视角须以独立材料复核。
\eventanchor{EQ-1784-TAIWAN-UPRISING-PRELUDE-DOCUMENTS}{围绕「台湾地方械斗、赋役与移民…}乾隆四十九年（1784年），围绕「台湾地方械斗、赋役与移民社会紧张为大规模反抗累积条件」的地方呈报、案件或赈务记录将社会反应带入官府文书链。核心取证：《清高宗实录》乾隆四十九年相关条；地方志、督抚奏折及地方档案。这一锚点记录的是文书链：它能证明呈报、上谕、部议或照会进入机构流程，不能由此推出实际执行。来源保留：此行仅确认相关文书链和可追查的行政动作；不得由其直接推断政策有效、地方服从、民众态度或单一因果。
\eventanchor{EQ-1785-TAIWAN-GARRISON-DOCUMENTS}{清廷围绕「清廷调整台湾驻防、渡…}乾隆五十年（1785年），清廷围绕「清廷调整台湾驻防、渡台与地方治安安排」处理盟旗、驻防、交通或地方呈报，相关边务文书进入中央决策链。核心取证：《清高宗实录》乾隆五十年相关条；军机处档、理藩院档或相关方略。这一锚点记录的是文书链：它能证明呈报、上谕、部议或照会进入机构流程，不能由此推出实际执行。来源保留：此行仅确认相关文书链和可追查的行政动作；不得由其直接推断政策有效、地方服从、民众态度或单一因果。
\eventanchor{EQ-1786-LIN-SHUANGWEN-UPRISING-DOCUMENTS}{清廷围绕「林爽文起事在台湾爆发…}乾隆五十一年（1786年），清廷围绕「林爽文起事在台湾爆发并迅速扩展」调遣军队、筹措饷械并处理战报，中央与地方的军政文书形成连续记录。核心取证：《清高宗实录》乾隆五十一年相关条；军机处档、战事方略及地方奏报。这一锚点记录的是文书链：它能证明呈报、上谕、部议或照会进入机构流程，不能由此推出实际执行。来源保留：此行仅确认相关文书链和可追查的行政动作；不得由其直接推断政策有效、地方服从、民众态度或单一因果。
\eventanchor{EQ-1786-FIRST-GURKHA-WAR-DOCUMENTS}{清廷围绕「廓尔喀军进入西藏边境…}乾隆五十一年（1786年），清廷围绕「廓尔喀军进入西藏边境，清廷开始处理高原军事危机」处理盟旗、驻防、交通或地方呈报，相关边务文书进入中央决策链。核心取证：《清高宗实录》乾隆五十一年相关条；军机处档、理藩院档或相关方略。这一锚点记录的是文书链：它能证明呈报、上谕、部议或照会进入机构流程，不能由此推出实际执行。来源保留：此行仅确认相关文书链和可追查的行政动作；不得由其直接推断政策有效、地方服从、民众态度或单一因果。
\eventanchor{EQ-1787-TAIWAN-CAMPAIGN-DOCUMENTS}{清廷围绕「福康安等率军赴台镇压…}乾隆五十二年（1787年），清廷围绕「福康安等率军赴台镇压林爽文起事，海峡运输成为关键」调遣军队、筹措饷械并处理战报，中央与地方的军政文书形成连续记录。核心取证：《清高宗实录》乾隆五十二年相关条；军机处档、战事方略及地方奏报。这一锚点记录的是文书链：它能证明呈报、上谕、部议或照会进入机构流程，不能由此推出实际执行。来源保留：此行仅确认相关文书链和可追查的行政动作；不得由其直接推断政策有效、地方服从、民众态度或单一因果。
\eventanchor{EQ-1788-LIN-SHUANGWEN-DEFEAT-DOCUMENTS}{清廷围绕「清军击败林爽文，台湾…}乾隆五十三年（1788年），清廷围绕「清军击败林爽文，台湾起事基本结束」调遣军队、筹措饷械并处理战报，中央与地方的军政文书形成连续记录。核心取证：《清高宗实录》乾隆五十三年相关条；军机处档、战事方略及地方奏报。这一锚点记录的是文书链：它能证明呈报、上谕、部议或照会进入机构流程，不能由此推出实际执行。来源保留：此行仅确认相关文书链和可追查的行政动作；不得由其直接推断政策有效、地方服从、民众态度或单一因果。
\eventanchor{EQ-1788-SECOND-GURKHA-WAR-DOCUMENTS}{清廷围绕「廓尔喀再次入侵西藏，…}乾隆五十三年（1788年），清廷围绕「廓尔喀再次入侵西藏，清军组织更大规模高原远征」调遣军队、筹措饷械并处理战报，中央与地方的军政文书形成连续记录。核心取证：《清高宗实录》乾隆五十三年相关条；军机处档、战事方略及地方奏报。这一锚点记录的是文书链：它能证明呈报、上谕、部议或照会进入机构流程，不能由此推出实际执行。来源保留：此行仅确认相关文书链和可追查的行政动作；不得由其直接推断政策有效、地方服从、民众态度或单一因果。
